% pdf-env-boxes.tex — עטיפת סביבות LaTeX קיימות עם tcolorbox, מותאם לצבעי HTML
\usepackage{amsthm}
\usepackage[most]{tcolorbox}
\tcbuselibrary{theorems,skins,breakable,hooks}
\usepackage{xcolor}

% ==== שמות בעברית ====
\providecommand{\definitionname}{הגדרה}
\providecommand{\examplename}{דוגמה}
\providecommand{\exercisename}{תרגיל}
\providecommand{\remarkname}{הערה}
\providecommand{\theoremname}{משפט}
\providecommand{\lemmaname}{למה}
\providecommand{\propositionname}{טענה}
\providecommand{\corollaryname}{מסקנה}
\renewcommand{\proofname}{הוכחה}

% ==== צבעים (בהתאם ל-CSS: rgba) ====
\definecolor{envDefinitionRGB}{RGB}{255,99,71}    % tomato
\definecolor{envExampleRGB}{RGB}{138,43,226}      % blueviolet
\definecolor{envExerciseRGB}{RGB}{30,144,255}     % dodgerblue
\definecolor{envRemarkRGB}{RGB}{0,0,0}            % black (עם שקיפות נמוכה)
\definecolor{envTheoremRGB}{RGB}{255,165,0}       % orange
\definecolor{envProofRGB}{RGB}{60,179,113}        % mediumseagreen

% ==== סגנון בסיסי לכל התיבות ====
\tcbset{
  envbox/.style={
    enhanced, breakable, sharp corners,
    boxrule=0.4pt,
    colback=white,
    fonttitle=\bfseries\small,
    coltitle=black,
    title filled,
    left=8pt,right=8pt,top=6pt,bottom=6pt,
    before skip=10pt, after skip=10pt
  },
}

\makeatletter
% parent counter: ב-book לפי chapter, אחרת לפי section
\@ifundefined{c@chapter}{\def\qnumwithin{section}}{\def\qnumwithin{chapter}}

\newcommand{\EnsureNewTheorem}[2]{%
  \@ifundefined{#1}{%
    \newtheorem{#1}{#2}[\qnumwithin]%
  }{}%
}
\makeatother

\AtBeginDocument{%
  \theoremstyle{definition}
  \EnsureNewTheorem{definition}{\definitionname}
  \EnsureNewTheorem{example}{\examplename}
  \EnsureNewTheorem{exercise}{\exercisename}

  \theoremstyle{remark}
  \EnsureNewTheorem{remark}{\remarkname}

  \theoremstyle{plain}
  \EnsureNewTheorem{theorem}{\theoremname}
  \EnsureNewTheorem{lemma}{\lemmaname}
  \EnsureNewTheorem{proposition}{\propositionname}
  \EnsureNewTheorem{corollary}{\corollaryname}

  % ==== עטיפות עם צבעי HTML ====
  \tcolorboxenvironment{definition}{envbox,
    colframe=envDefinitionRGB,
    colbacktitle=envDefinitionRGB,
    colback=envDefinitionRGB!10,
    borderline south={0.4pt}{0pt}{envDefinitionRGB},
    title={\definitionname}}

  \tcolorboxenvironment{example}{envbox,
    colframe=envExampleRGB,
    colbacktitle=envExampleRGB,
    colback=envExampleRGB!10,
    borderline south={0.4pt}{0pt}{envExampleRGB},
    title={\examplename}}

  \tcolorboxenvironment{exercise}{envbox,
    colframe=envExerciseRGB,
    colbacktitle=envExerciseRGB,
    colback=envExerciseRGB!10,
    borderline south={0.4pt}{0pt}{envExerciseRGB},
    title={\exercisename}}

  \tcolorboxenvironment{remark}{envbox,
    colframe=envRemarkRGB,
    colbacktitle=envRemarkRGB,
    colback=envRemarkRGB!10,
    borderline south={0.4pt}{0pt}{envRemarkRGB},
    title={\remarkname}}

  \tcolorboxenvironment{theorem}{envbox,
    colframe=envTheoremRGB,
    colbacktitle=envTheoremRGB,
    colback=envTheoremRGB!10,
    borderline south={0.4pt}{0pt}{envTheoremRGB},
    title={\theoremname}}

  \tcolorboxenvironment{lemma}{envbox,
    colframe=envTheoremRGB,
    colbacktitle=envTheoremRGB,
    colback=envTheoremRGB!10,
    borderline south={0.4pt}{0pt}{envTheoremRGB},
    title={\lemmaname}}

  \tcolorboxenvironment{proposition}{envbox,
    colframe=envTheoremRGB,
    colbacktitle=envTheoremRGB,
    colback=envTheoremRGB!10,
    borderline south={0.4pt}{0pt}{envTheoremRGB},
    title={\propositionname}}

  \tcolorboxenvironment{corollary}{envbox,
    colframe=envTheoremRGB,
    colbacktitle=envTheoremRGB,
    colback=envTheoremRGB!10,
    borderline south={0.4pt}{0pt}{envTheoremRGB},
    title={\corollaryname}}

  \tcolorboxenvironment{proof}{envbox,
    colframe=envProofRGB,
    colbacktitle=envProofRGB,
    colback=envProofRGB!10,
    borderline south={0.4pt}{0pt}{envProofRGB},
    title={\proofname}}
}
